% \iffalse meta-comment
% LTeX: language=de-DE
%
%  TUDCD-Script -- Corporate Design of Technische Universität Dresden
% ----------------------------------------------------------------------------
%
%  Copyright (C) TU Dresden <tudlatex@tu-dresden.de>, 2025
%
% ----------------------------------------------------------------------------
%
%  This work may be distributed and/or modified under the conditions of the
%  LaTeX Project Public License, either version 1.3c of this license or
%  any later version. The latest version of this license is in
%    http://www.latex-project.org/lppl.txt
%  and version 1.3c or later is part of all distributions of
%  LaTeX version 2008-05-04 or later.
%
%  This work has the LPPL maintenance status "maintained".
%
%  The current maintainer and author of this work is Yoshi Diepelt.
%
% ----------------------------------------------------------------------------
% \fi
% \selectlanguage{ngerman}
%
% \subsection{\texttt{test-compile}: Übersetzungstest für Fehlerhafte und korrekte Übersetzungen}
%
% Diese Folge an Test ist genau dann erfolgreich, wenn die Kompilation
% ohne Fehler abbricht und eine PDF generiert.
%
% Mit dem folgenden Guard wird die bereitgestellte Testsuite von \dpkg{l3build}
% in alle Tests dieser Suite eingebunden.
% Da die Batchdatei zum erstellen des Tests zerstückelt geschrieben wird,
% müssen Guards um die Präambel der Tests geschrieben werden.
% Beim Entpacken werden die Testdokumente innerhalb einer \dcode{\textbackslash ifx}-Abfrage übersprungen.
% Damit darin enthaltene \dmacro{\iftutex}-Abfragen korrekt verschachtelt werden, muss \dmacro{\iftutex}
% beim Entpacken als Bedingung bekannt sein.
%    \begin{macrocode}
%<*ins>
\input docstrip.tex
\def\tudcdbatchfilemode{\relax}
\let\iftutex\iffalse
%</ins>
%<*ins>
\ifx\tudcdbatchfilemode\undefined%
%</ins>
%<*preamble>
\input{regression-test.tex}
%</preamble>
%<*ins>
\fi
%</ins>
%    \end{macrocode}
%
% \begin{unittest}{dev-version}{Kompilation der Artikelklasse}
%
% In dem ersten Test wird eine erfolgreiche Kompilation der Artikelklasse
% getestet.
%
%    \begin{macrocode}
%<*ins>
\generate{\file{\jobname-0001.lvt}{\from{\jobname.dtx}{0001,preamble}}}
%</ins>
%    \end{macrocode}
% Und dies ohne jegliche zusätzliche Pakete.
%    \begin{macrocode}
%<*ins>
\ifx\tudcdbatchfilemode\undefined%
%</ins>
%<*0001>
\documentclass{tudcdartcl}

\begin{document}
Hallo Welt!
\end{document}
%</0001>
%<*ins>
\fi
%</ins>
%    \end{macrocode}
%
% \end{unittest}
%
% \begin{unittest}{dev-version}{Kompilation der Reportklasse}
%
% In diesem Test wird eine erfolgreiche Kompilation der Reportklasse
% getestet.
%    \begin{macrocode}
%<*ins>
\generate{\file{\jobname-0002.lvt}{\from{\jobname.dtx}{0002,preamble}}}
%</ins>
%    \end{macrocode}
% Und dies ohne jegliche zusätzliche Pakete.
%    \begin{macrocode}
%<*ins>
\ifx\tudcdbatchfilemode\undefined%
%</ins>
%<*0002>
\documentclass{tudcdreprt}

\begin{document}
Hallo Welt!
\end{document}
%</0002>
%<*ins>
\fi
%</ins>
%    \end{macrocode}
% \end{unittest}
%
% \begin{unittest}{dev-version}{Kompilation der Buchklasse}
%
% In diesem Test wird eine erfolgreiche Kompilation der Buchklasse
% getestet.
%    \begin{macrocode}
%<*ins>
\generate{\file{\jobname-0003.lvt}{\from{\jobname.dtx}{0003,preamble}}}
%</ins>
%    \end{macrocode}
% Und dies ohne jegliche zusätzliche Pakete.
%    \begin{macrocode}
%<*ins>
\ifx\tudcdbatchfilemode\undefined%
%</ins>
%<*0003>
\documentclass{tudcdbook}

\begin{document}
Hallo Welt!
\end{document}
%</0003>
%<*ins>
\fi
%</ins>
%    \end{macrocode}
%
% \end{unittest}
%
% \begin{unittest}{dev-version}{Kompilation der Briefklasse}
%
% In diesem Test wird eine erfolgreiche Kompilation der Briefklasse
% getestet.
%    \begin{macrocode}
%<*ins>
\generate{\file{\jobname-0004.lvt}{\from{\jobname.dtx}{0004,preamble}}}
%</ins>
%    \end{macrocode}
% Und dies ohne jegliche zusätzliche Pakete.
%    \begin{macrocode}
%<*ins>
\ifx\tudcdbatchfilemode\undefined%
%</ins>
%<*0004>
\documentclass{tudcdlttr}

\begin{document}
Hallo Welt!
\end{document}
%</0004>
%<*ins>
\fi
%</ins>
%    \end{macrocode}
%
% \end{unittest}
%
% \begin{unittest}{dev-version}{Kompilation der Beamerklasse}
%
% In diesem Test wird eine erfolgreiche Kompilation der Beamerklasse
% getestet.
%    \begin{macrocode}
%<*ins>
\generate{\file{\jobname-0005.lvt}{\from{\jobname.dtx}{0005,preamble}}}
%</ins>
%    \end{macrocode}
% Und dies ohne jegliche zusätzliche Pakete.
%    \begin{macrocode}
%<*ins>
\ifx\tudcdbatchfilemode\undefined%
%</ins>
%<*0005>
\documentclass{beamer}

\usetheme{tudcd}

\begin{document}
\begin{frame}
  \frametitle{Beamertest}
  \framesubtitle{Beamertest subtitle}
  Hello World!
\end{frame}
\end{document}
%</0005>
%<*ins>
\fi
%</ins>
%    \end{macrocode}
%
% \end{unittest}
%
% \begin{unittest}{dev-version}{Kompilation der Klasse für Abschlussarbeiten}
%
% In diesem Test wird eine erfolgreiche Kompilation der Klasse \dcode{tudcdthesis} ohne zusätzliche Pakete getestet.
%    \begin{macrocode}
%<*ins>
\generate{\file{\jobname-0006.lvt}{\from{\jobname.dtx}{0006,preamble}}}
%</ins>
%    \end{macrocode}
%
%    \begin{macrocode}
%<*ins>
\ifx\tudcdbatchfilemode\undefined%
%</ins>
%<*0006>
\documentclass{tudcdthesis}
\begin{document}
Hallo Welt!
\end{document}
%</0006>
%<*ins>
\fi
%</ins>
%    \end{macrocode}
%
% \end{unittest}
%
% \begin{unittest}{dev-version}{Alle Elemente der Klasse für Abschlussarbeiten}
%
% Dieser Test setzt alle Eingabefelder, die Titelseite, die Zusammenfassung und die Erklärungen im Druckmodus.
%    \begin{macrocode}
%<*ins>
\generate{\file{\jobname-0007.lvt}{\from{\jobname.dtx}{0007,preamble}}}
%</ins>
%    \end{macrocode}
%
%    \begin{macrocode}
%<*ins>
\ifx\tudcdbatchfilemode\undefined%
%</ins>
%<*0007>
\documentclass[output=print,abstract={section,totoc},declaration=chapter]{tudcdthesis}
\usepackage[english,ngerman]{babel}
\usepackage{hyperref}
\faculty{Fakultät Maschinenwesen}
\institute{Institut für Festkörpermechanik}
\chair{Professur für Dynamik und Mechanismentechnik}
\thesis{diploma}
\graduation[Dipl.-Ing.]{Diplomingenieur}
\title{Über das ordentliche Setzen einer Abschlussarbeit}
\subtitle{im Corporate Design der TU~Dresden}
\author{Hans Wirdingenieur\matriculationnumber{1234567}\emailaddress{hans@example.org}}
\dateofbirth{20.02.2002}
\placeofbirth{Dresden}
\supervisor{Paul Müller \and Paula Schmidt}
\referee{Prof. A \and Prof. B}
\professor{Prof. Max Mustermann}
\defensedate{1.1.2027}
\date{24. Dezember 2026}
\begin{document}
\frontmatter
\maketitle
\begin{abstract}
Kurzfassung.
\nextabstract[english]
Abstract.
\end{abstract}
\begin{declarations}[place=Leipzig]
\confirmation*
\blocking[company=Muster GmbH]
\nextdeclaration{Eigene Erklärung}{Text}
\end{declarations}
\tableofcontents
\mainmatter
\chapter{Einleitung}
\section{Abschnitt}
Text.
\end{document}
%</0007>
%<*ins>
\fi
%</ins>
%    \end{macrocode}
%
% \end{unittest}
%
% \begin{unittest}{dev-version}{Abwärtskompatibilität: Beispiel \texttt{thesis.tex} aus \texttt{tudscr}}
%
% Das Beispiel einer Masterarbeit zweier Autoren aus \dpkg{tudscr}, in welchem lediglich die Klasse ersetzt wurde.
%    \begin{macrocode}
%<*ins>
\generate{\file{\jobname-0008.lvt}{\from{\jobname.dtx}{0008,preamble}}}
%</ins>
%    \end{macrocode}
%
%    \begin{macrocode}
%<*ins>
\ifx\tudcdbatchfilemode\undefined%
%</ins>
%<*0008>
\documentclass[ngerman,BCOR=4mm]{tudcdthesis}
\iftutex
  \usepackage{fontspec}
\else
  \usepackage[T1]{fontenc}
  \usepackage[ngerman=ngerman-x-latest]{hyphsubst}
\fi
\usepackage{babel}
\usepackage{isodate}
\begin{document}
\faculty{Juristische Fakultät}
\department{Fachrichtung Strafrecht}
\institute{Institut für Kriminologie}
\chair{Lehrstuhl für Kriminalprognose}
\date{18.02.2015}
\title{%
  Entwicklung eines optimalen Verfahrens zur Eroberung des
  Geldspeichers in Entenhausen
}
\subject{master}
\graduation[M.Sc.]{Master of Science}
\author{%
  Mickey Mouse%
  \matriculationnumber{12345678}%
  \dateofbirth{2.1.1990}%
  \placeofbirth{Dresden}%
\and%
  Donald Duck%
  \matriculationnumber{87654321}%
  \dateofbirth{1.2.1990}%
  \placeofbirth{Berlin}%
}
\matriculationyear{2010}
\supervisor{Dagobert Duck \and Mac Moneysac}
\professor{Prof. Dr. Kater Karlo}
\maketitle
\end{document}
%</0008>
%<*ins>
\fi
%</ins>
%    \end{macrocode}
%
% \end{unittest}
%
% \begin{unittest}{dev-version}{Abwärtskompatibilität: Beispiel \texttt{dissertation.tex} aus \texttt{tudscr}}
%
% Das Beispiel einer Dissertation aus \dpkg{tudscr}, ergänzt um einen eigenen Bezeichner mit \dmacro{\getfield}, \dmacro{\TUDoptions} mit einer Layoutoption von \dpkg{tudscr} sowie \dmacro{\declaration}.
%    \begin{macrocode}
%<*ins>
\generate{\file{\jobname-0009.lvt}{\from{\jobname.dtx}{0009,preamble}}}
%</ins>
%    \end{macrocode}
%
%    \begin{macrocode}
%<*ins>
\ifx\tudcdbatchfilemode\undefined%
%</ins>
%<*0009>
\documentclass[%
  english,ngerman,%
  BCOR=6mm,cdgeometry=no,DIV=13%
]{tudcdthesis}
\iftutex
  \usepackage{fontspec}
\else
  \usepackage[T1]{fontenc}
  \usepackage[ngerman=ngerman-x-latest]{hyphsubst}
\fi
\usepackage{babel}
\usepackage{isodate}
\usepackage{blindtext}
\renewcaptionname{ngerman}{\confirmationtext}{%
  Ich versichere, die Arbeit \glqq\getfield{title}\grqq{} selbstständig verfasst zu haben.}
\begin{document}
\faculty{Juristische Fakultät}
\department{Fachrichtung Strafrecht}
\institute{Institut für Kriminologie}
\chair{Lehrstuhl für Kriminalprognose}
\date{18.02.2015}
\author{Mickey Mouse}
\title{%
  Entwicklung eines optimalen Verfahrens zur Eroberung des
  Geldspeichers in Entenhausen
}
\thesis{diss}
\graduation[Dr. jur.]{Doktor der Rechte}
\dateofbirth{2.1.1990}
\placeofbirth{Dresden}
\defensedate{15.03.2015}
\referee{Dagobert Duck \and Mac Moneysac}
\maketitle
\TUDoption{abstract}{section,multiple}
\begin{abstract}[pagestyle=empty.tudheadings]
\blindtext
\nextabstract[english]
\blindtext
\end{abstract}
\TUDoptions{cd=true,parskip=half}
\declaration[company=Entenhausen AG]
\end{document}
%</0009>
%<*ins>
\fi
%</ins>
%    \end{macrocode}
%
% \end{unittest}
%
%    \begin{macrocode}
%<*ins>
\endbatchfile%
%</ins>
%    \end{macrocode}
%
